\documentclass[aspectratio=169]{beamer}

% Tema UFS --- Universidade Federal de Sergipe
\usetheme{UFS}

% Pacotes base
\usepackage[utf8]{inputenc}
\usepackage[portuguese]{babel}
\usepackage{amsmath,amssymb}
\usepackage{graphicx}
\usepackage{booktabs}
\usepackage{xcolor}

% TikZ e bibliotecas
\usepackage{tikz}
\usetikzlibrary{shapes.geometric, arrows.meta, positioning, matrix, fit,
  backgrounds, decorations.pathreplacing, shadows, patterns, calc}

% Listagens --- paleta VS Code Light+
\usepackage{listings}
\definecolor{bgcolor}{HTML}{FFFFFF}
\definecolor{linecolor}{HTML}{DDDDDD}
\definecolor{numcolor}{HTML}{999999}
\definecolor{kwcolor}{HTML}{0000FF}
\definecolor{strcolor}{HTML}{A31515}
\definecolor{cmtcolor}{HTML}{008000}

\lstdefinestyle{ufsStyle}{
  backgroundcolor=\color{bgcolor},
  basicstyle=\ttfamily\footnotesize\color{black},
  keywordstyle=\color{kwcolor}\bfseries,
  stringstyle=\color{strcolor},
  commentstyle=\color{cmtcolor}\itshape,
  numberstyle=\tiny\color{numcolor},
  numbers=left, numbersep=12pt,
  xleftmargin=20pt, framexleftmargin=20pt,
  breaklines=true, tabsize=4,
  keepspaces=true, showspaces=false,
  showstringspaces=false, showtabs=false,
  frame=single, rulecolor=\color{linecolor},
  framesep=4pt, captionpos=b, upquote=true,
  columns=fullflexible,
}
\lstset{style=ufsStyle, language=C}

\lstdefinestyle{tinyC}{
  style=ufsStyle, language=C, numbers=none,
  basicstyle=\ttfamily\scriptsize\color{black},
}
\lstdefinestyle{shellStyle}{
  style=ufsStyle, language=bash, numbers=none,
  basicstyle=\ttfamily\scriptsize\color{black},
}
\lstdefinestyle{makeStyle}{
  style=ufsStyle, language=make, numbers=none,
  basicstyle=\ttfamily\scriptsize\color{black},
}

% --- Estilos TikZ reutilizados ---
\tikzset{
  mod/.style={rectangle, draw=UFSBlue, fill=UFSBlue!8, rounded corners,
    minimum width=2.6cm, minimum height=0.85cm, font=\small\ttfamily,
    align=center},
  modh/.style={rectangle, draw=UFSYellow!80!black, fill=UFSYellow!20,
    rounded corners, minimum width=2.2cm, minimum height=0.7cm,
    font=\small\ttfamily, align=center},
  modr/.style={rectangle, draw=UFSRed, fill=UFSRed!8, rounded corners,
    minimum width=2.6cm, minimum height=0.85cm, font=\small\ttfamily,
    align=center},
  obj/.style={rectangle, draw=UFSGray, fill=black!5, minimum width=1.6cm,
    minimum height=0.7cm, font=\small\ttfamily},
  rot/.style={font=\scriptsize\ttfamily, text=UFSGray},
  leg/.style={font=\scriptsize, text=UFSGray},
  seta/.style={->, thick, UFSBlue},
  setar/.style={->, thick, UFSRed},
  setag/.style={->, thick, UFSGray},
  caixa/.style={rectangle, draw=UFSGray, rounded corners, fill=black!3,
    font=\scriptsize, align=center, inner sep=4pt},
}

% Referencia da fonte no rodape do frame
\newcommand{\fonte}[1]{\vfill\hfill{\scriptsize\textcolor{UFSGray}{#1}}}

% --- Metadados ---
\title{Padrões de codificação}
\subtitle{E a prática de modularizar seguindo um padrão}
\author{Prof.\ André Yoshiaki Kashiwabara}
\institute{DCOMP -- Universidade Federal de Sergipe\\
\texttt{andreyoshiaki@dcomp.ufs.br}}
\date{}

\begin{document}

\begin{frame}[plain]
  \titlepage
\end{frame}

\begin{frame}{O que você vai aprender hoje}
  \begin{center}
  \begin{tikzpicture}[
    item/.style={rectangle, draw=UFSBlue, fill=UFSBlue!10, rounded corners,
      minimum width=10.4cm, minimum height=0.72cm, font=\small}
  ]
    \node[item] (t1) at (0, 0)    {Por que uma equipe adota um padrão de codificação};
    \node[item] (t2) at (0,-0.95) {Nomes e formatação: o que o padrão fixa};
    \node[item] (t3) at (0,-1.90) {Comentários que valem o espaço que ocupam};
    \node[item] (t4) at (0,-2.85) {\texttt{static}, \texttt{const}, constantes com nome, \texttt{-Werror}};
    \node[item] (t5) at (0,-3.80) {Aplicar tudo isso ao módulo que vamos encaixar no projeto};
    \draw[seta] (t1) -- (t2);
    \draw[seta] (t2) -- (t3);
    \draw[seta] (t3) -- (t4);
    \draw[seta] (t4) -- (t5);
  \end{tikzpicture}
  \end{center}
\end{frame}

\begin{frame}{O problema: um exemplo de código ruim que funciona}
  \begin{columns}[T]
    \begin{column}{0.5\textwidth}
      \textbf{Onde estamos}\\[0.3em]
      Na aula passada o \texttt{notas.c} virou quatro arquivos:
      \texttt{matriz}, \texttt{estat}, \texttt{turma} e o \texttt{main}.

      \medskip
      Hoje um colega manda o \texttt{ranking.c}: classifica os alunos pela
      média. Ele jura que funciona --- e funciona.
    \end{column}
    \begin{column}{0.5\textwidth}
      \textbf{O que vamos descobrir}\\[0.3em]
      \textcolor{UFSRed}{Não tem \texttt{.h}:} para chamá-lo, o \texttt{main.c}
      copiou as assinaturas na mão.

      \medskip
      \textcolor{UFSRed}{Guarda estado em variáveis globais.}

      \medskip
      \textcolor{UFSRed}{Ninguém consegue ler} sem parar para decifrar.
    \end{column}
  \end{columns}
  \begin{block}{A pergunta de hoje}
    O que separa um código que \emph{funciona} de um código que a turma toda
    consegue manter?
  \end{block}
\end{frame}

% ==================================================================
\section{Um exemplo de código ruim}
% ==================================================================

\begin{frame}[fragile]{Exemplo de código ruim: o estado}
\begin{lstlisting}[style=tinyC]
double R[50];
int Q[50];
int N;

void troca(int a,int b){int t=Q[a];Q[a]=Q[b];Q[b]=t;}

int maior(double x,double y){ return x>y; }
\end{lstlisting}
  \begin{block}{}
    \small Três variáveis globais e duas funções auxiliares que qualquer
    outro arquivo do projeto enxerga.
  \end{block}
\end{frame}

\begin{frame}[fragile]{Exemplo de código ruim: o cálculo}
\begin{lstlisting}[style=tinyC]
void calcula(Turma *t, int flag)
{
  double tmp;
  N = t->nalunos;
  for(int i=0;i<N;i++){
  R[i]=media(t->notas[i],t->navaliacoes);
    Q[i]=i;                 /* guarda o indice i em Q[i] */
  }
	for(int i=0;i<N;i++)
		for(int j=0;j<N-1;j++)
			if(maior(R[Q[j+1]],R[Q[j]])) troca(j,j+1);
}
\end{lstlisting}
\end{frame}

\begin{frame}[fragile]{...e como o \texttt{main.c} conseguiu chamá-lo}
\begin{lstlisting}[style=tinyC]
void mostra()
{
    printf("%-5s %-8s %8s %s\n","pos","aluno","media","situacao");
  for(int i=0;i<N;i++){
    printf("%-5d %-8d %8.2f %s\n", i+1, Q[i], R[Q[i]],
           R[Q[i]]>=7.0?"aprovado":"reprovado");
  }
}
\end{lstlisting}
\begin{lstlisting}[style=tinyC]
/* no main.c: o ranking nao tem header,
   copiei as assinaturas na mao */
void calcula(Turma *t, int flag);
void mostra();
\end{lstlisting}
\end{frame}

\begin{frame}{O que vai acontecer?}
  \begin{block}{Antes de rodar, decida}
    \begin{enumerate}
      \item O projeto compila com \texttt{-Wall -Wextra}? Com quantos avisos?
      \item A saída do ranking sai certa para os três alunos do
        \texttt{main.c}?
      \item E se a turma tiver 60 alunos em vez de 3 --- o que acontece?
      \item Amanhã o \texttt{turma.c} ganha uma função chamada
        \texttt{troca}. O que acontece na ligação?
    \end{enumerate}
  \end{block}
  \begin{exampleblock}{}
    Anote suas quatro respostas antes de executar \texttt{make}.
  \end{exampleblock}
\end{frame}

\begin{frame}[fragile]{O que aconteceu}
\begin{lstlisting}[style=shellStyle]
$ make
ranking.c:16:10: warning: unused variable 'tmp' [-Wunused-variable]
ranking.c:14:28: warning: unused parameter 'flag' [-Wunused-parameter]
$ ./notas
pos   aluno       media situacao
1     2            9.25 aprovado
2     0            7.62 aprovado
3     1            5.75 reprovado
\end{lstlisting}
  \begin{block}{}
    \small Perguntas 1 e 2 respondidas: compila com dois avisos, roda e
    acerta. As perguntas 3 e 4 a execução não mostrou --- e são as que
    interessam.
  \end{block}
\end{frame}

\begin{frame}[fragile]{Pergunta 3: o limite que ninguém escreveu}
  \begin{center}
  \begin{tikzpicture}[scale=1]
    \draw[fill=UFSBlue!10, draw=UFSBlue] (0,0) rectangle (5.6,0.65);
    \node[font=\scriptsize\ttfamily] at (2.8,0.32) {R[0] ... R[49]};
    \draw[fill=black!4, draw=UFSGray, dashed] (5.6,0) rectangle (8.6,0.65);
    \node[font=\scriptsize] at (7.1,0.32) {o que vier depois};
    \draw[setar] (6.6,1.35) -- (6.6,0.7);
    \node[leg, align=center] at (6.6,1.62)
      {a turma de 60 escreve \texttt{R[50]} ... \texttt{R[59]} aqui};
  \end{tikzpicture}
  \end{center}
  \begin{columns}[T]
    \begin{column}{0.5\textwidth}
      \small \textcolor{UFSRed}{\textbf{Em uma máquina}} depois de
      \texttt{R} vinha \texttt{Q}: os índices dos alunos viram lixo e as
      últimas colocações saem com média \texttt{0.00}.
    \end{column}
    \begin{column}{0.5\textwidth}
      \small \textcolor{UFSRed}{\textbf{Em outra}} depois de \texttt{R}
      vinha \texttt{N}: a contagem vira 50 e o programa imprime 50 linhas
      para uma turma de 60.
    \end{column}
  \end{columns}
  \begin{block}{}
    \small Não existe ``o que acontece'': existe comportamento indefinido.
    O programa não quebrou --- ele mentiu, e mentiu diferente em cada
    máquina. É o desfecho normal quando um limite fica implícito num número
    solto.
  \end{block}
  \fonte{ISO/IEC 9899:2018, §J.2 (comportamento indefinido)}
\end{frame}

\begin{frame}[fragile]{Pergunta 4: dois \texttt{troca} no mesmo programa}
\begin{lstlisting}[style=shellStyle]
$ gcc main.o turma.o matriz.o estat.o ranking.o outro.o -o notas
duplicate symbol '_troca' in: outro.o / ranking.o
ld: 1 duplicate symbols          # no Linux: "multiple definition of 'troca'"
\end{lstlisting}
  \begin{columns}[T]
    \begin{column}{0.5\textwidth}
      \small \textbf{O que aconteceu}\\
      Nada está errado no \texttt{outro.c}: ele só precisou de uma função
      \texttt{troca}. Quem reservou um nome comum para o projeto inteiro foi
      o \texttt{ranking.c}.
    \end{column}
    \begin{column}{0.5\textwidth}
      \small \textbf{O custo}\\
      O erro aparece na \emph{ligação}, longe do arquivo que o causou, e
      recai sobre quem chegou depois. Quanto maior o projeto, mais caro.
    \end{column}
  \end{columns}
  \fonte{ISO/IEC 9899:2018, §6.2.2 (ligação externa e interna)}
\end{frame}

\begin{frame}{Funciona não é o mesmo que está pronto}
  \begin{columns}[T]
    \begin{column}{0.52\textwidth}
      \textbf{Quem lê esse código depois}\\[0.3em]
      Você daqui a três meses. O colega de equipe. Quem for consertar o bug
      às onze da noite.

      \medskip
      Um trecho é escrito uma vez e lido dezenas --- o tempo de leitura é o
      custo que importa.
    \end{column}
    \begin{column}{0.48\textwidth}
      \textbf{O que um padrão de codificação é}\\[0.3em]
      Um conjunto de decisões tomadas \emph{uma vez}, por escrito, para não
      serem rediscutidas a cada arquivo: nomes, formatação, comentários,
      organização dos arquivos.

      \medskip
      Não é gosto pessoal: é acordo.
    \end{column}
  \end{columns}
  \fonte{Kernighan \& Pike (1999), \emph{The Practice of Programming}, cap.~1}
\end{frame}

% ==================================================================
\section{Nomes}
% ==================================================================

\begin{frame}{Nomes dizem o que a coisa é}
  \begin{center}
  \small
  \begin{tabular}{l l l}
    \toprule
    \textbf{No código ruim} & \textbf{O que é de verdade} &
      \textbf{Nome} \\
    \midrule
    \texttt{R[i]} & média do aluno \texttt{i} & \texttt{media\_aluno[i]} \\
    \texttt{Q[i]} & quem está na $i$-ésima colocação & \texttt{posicao[i]} \\
    \texttt{N} & quantos alunos entraram no ranking & \texttt{quantidade} \\
    \texttt{calcula} & calcula \emph{o quê}? & \texttt{ranking\_calcula} \\
    \texttt{mostra} & mostra \emph{o quê}? & \texttt{ranking\_imprime} \\
    \texttt{flag} & nada: ninguém usa & --- \\
    \bottomrule
  \end{tabular}
  \end{center}
  \begin{block}{}
    \small Na Tarefa 2, quando o estado sai das globais e passa a viver numa
    \texttt{struct}, \texttt{ranking\_calcula} vira \texttt{ranking\_cria}:
    aí a função passa mesmo a \emph{criar} alguma coisa.
  \end{block}
  \fonte{Kernighan \& Pike (1999), §1.1 ``Names''; McConnell (2004), cap.~11}
\end{frame}

\begin{frame}{A regra do escopo, e o prefixo do módulo}
  \begin{columns}[T]
    \begin{column}{0.5\textwidth}
      \textbf{Escopo curto, nome curto}\\[0.3em]
      \texttt{i}, \texttt{j}, \texttt{n} em um laço de três linhas estão
      certos --- todo mundo lê sem esforço.

      \medskip
      \textbf{Escopo longo, nome descritivo}\\[0.3em]
      O que é global, ou faz parte da interface, paga por um nome inteiro.
    \end{column}
    \begin{column}{0.5\textwidth}
      \textbf{Prefixo de módulo}\\[0.3em]
      O projeto já tem a convenção:\\[0.3em]
      \texttt{turma\_cria}, \texttt{turma\_libera},
      \texttt{turma\_relatorio}

      \medskip
      O módulo novo entra igual:\\[0.3em]
      \texttt{ranking\_cria}, \texttt{ranking\_libera},
      \texttt{ranking\_imprime}

      \medskip
      C não tem \emph{namespace}: o prefixo é o nosso.
    \end{column}
  \end{columns}
  \fonte{Kernighan \& Pike (1999), §1.1; GNU Coding Standards, §5.4}
\end{frame}

% ==================================================================
\section{Formatação}
% ==================================================================

\begin{frame}[fragile]{O mesmo laço, depois do \texttt{clang-format}}
  \begin{columns}[T]
    \begin{column}{0.48\textwidth}
      \textcolor{UFSRed}{\textbf{Código ruim}}
\begin{lstlisting}[style=tinyC]
  for(int i=0;i<N;i++){
  R[i]=media(t->notas[i],
    t->navaliacoes);
    Q[i]=i;
  }
	for(int i=0;i<N;i++)
		for(int j=0;j<N-1;j++)
			if(maior(R[Q[j+1]],R[Q[j]]))
        troca(j,j+1);
\end{lstlisting}
    \end{column}
    \begin{column}{0.48\textwidth}
      \textcolor{UFSBlue}{\textbf{Formatado}}
\begin{lstlisting}[style=tinyC]
for (int i = 0; i < N; i++) {
    R[i] = media(t->notas[i],
                 t->navaliacoes);
    Q[i] = i;
}
for (int i = 0; i < N; i++)
    for (int j = 0; j < N - 1; j++)
        if (maior(R[Q[j+1]], R[Q[j]]))
            troca(j, j + 1);
\end{lstlisting}
    \end{column}
  \end{columns}
  \begin{block}{}
    \small Byte por byte, o mesmo programa: só mudou o tempo que você gasta
    para ter certeza disso. Repare no que \emph{não} mudou --- \texttt{R},
    \texttt{Q}, \texttt{N} e \texttt{maior} continuam lá.
  \end{block}
\end{frame}

\begin{frame}{O que um padrão de formatação fixa}
  \begin{columns}[T]
    \begin{column}{0.5\textwidth}
      \begin{itemize}
        \item indentação: 4 espaços, nunca tabulação (ou o contrário --- mas
          um só);
        \item largura máxima de linha: 80 colunas;
        \item onde a chave de abertura fica: função em linha nova, bloco na
          mesma linha;
      \end{itemize}
    \end{column}
    \begin{column}{0.5\textwidth}
      \begin{itemize}
        \item espaço depois de \texttt{if}, \texttt{for}, \texttt{while};
          nenhum antes do parêntese da chamada;
        \item uma declaração por linha;
        \item ordem dentro do \texttt{.c}: \texttt{include}s, constantes,
          tipos, funções \texttt{static}, funções da interface.
      \end{itemize}
    \end{column}
  \end{columns}
  \fonte{Linux kernel coding style, §1--§3; GNU Coding Standards, §5.1}
\end{frame}

\begin{frame}{Consistência vale mais do que preferência}
  \begin{block}{Por que não adianta discutir qual é o melhor}
    Chave na mesma linha ou na linha de baixo? Dois espaços ou quatro? Não
    há resposta certa --- há projetos que escolheram. O ganho vem de
    \emph{todo o arquivo} seguir a mesma escolha, porque aí a formatação
    deixa de chamar atenção e você lê o que o código faz.
  \end{block}
  \begin{exampleblock}{Regra prática}
    Entrou em código alheio, siga o estilo dele. Um arquivo com dois estilos
    é pior do que um arquivo com o estilo que você não gosta.
  \end{exampleblock}
  \fonte{Kernighan \& Pike (1999), §1.3 ``Consistency and Idioms''}
\end{frame}

% ==================================================================
\section{Comentários}
% ==================================================================

\begin{frame}[fragile]{Três comentários e o que cada um vale}
\begin{lstlisting}[style=tinyC]
Q[i] = i;                    /* guarda o indice i em Q[i] */
\end{lstlisting}
  \small\textcolor{UFSRed}{Repete o código.} Se a linha mudar, ele mente.
  \medskip
\begin{lstlisting}[style=tinyC]
/* ordena por media decrescente (usa selecao) */
... insercao ...
\end{lstlisting}
  \small\textcolor{UFSRed}{Desatualizado.} Pior que nenhum: o leitor acredita
  nele.
  \medskip
\begin{lstlisting}[style=tinyC]
/* Insercao, e nao um algoritmo melhor, porque
   uma turma cabe em uma sala. */
\end{lstlisting}
  \small\textcolor{UFSBlue}{Explica a decisão} --- o que o código não
  consegue dizer sozinho.
  \fonte{McConnell (2004), cap.~32; Kernighan \& Pike (1999), §1.6}
\end{frame}

\begin{frame}{Onde cada comentário mora}
  \begin{columns}[T]
    \begin{column}{0.5\textwidth}
      \textbf{No \texttt{.h}: o contrato}\\[0.3em]
      Para quem vai \emph{usar} sem abrir o \texttt{.c}:
      \begin{itemize}
        \item o que a função faz;
        \item o que devolve, inclusive no erro;
        \item de quem é a responsabilidade de liberar.
      \end{itemize}
    \end{column}
    \begin{column}{0.5\textwidth}
      \textbf{No \texttt{.c}: o porquê}\\[0.3em]
      Para quem vai \emph{manter}:
      \begin{itemize}
        \item por que este algoritmo e não outro;
        \item que caso estranho aquela linha resolve;
        \item que suposição está sendo feita.
      \end{itemize}
    \end{column}
  \end{columns}
  \begin{block}{}
    \small O comentário que descreve \emph{como} costuma ser sintoma: quase
    sempre dá para trocá-lo por um nome melhor ou por uma função à parte.
  \end{block}
  \fonte{McConnell (2004), cap.~32; GNU Coding Standards, §5.2}
\end{frame}

\begin{frame}[fragile]{Sua vez: dois minutos, em dupla}
\begin{lstlisting}[style=tinyC]
int p(double *v, int n)
{
    double s = 0;              /* soma */
    for (int i = 0; i < n; i++)
        s += v[i];
    if (s / n >= 7.0)          /* compara com 7.0 */
        return 1;
    return 0;
}
\end{lstlisting}
  \begin{block}{}
    Esta função já está formatada no padrão --- e mesmo assim três coisas
    dela quebram o que combinamos até agora. Quais? Escreva o nome que você
    daria à função, aos parâmetros e ao \texttt{7.0}.
  \end{block}
\end{frame}

\begin{frame}{O que dava para melhorar}
  \begin{columns}[T]
    \begin{column}{0.5\textwidth}
      \small
      \textbf{1. Nomes de escopo curto num escopo longo}\\
      \texttt{p} e \texttt{v} estão na interface:
      \texttt{aprovado\_pela\_media(const double *notas, int quantidade)}.

      \medskip
      \textbf{2. Comentários que repetem o código}\\
      ``soma'' e ``compara com 7.0'' somem. Falta o que o código não diz: o
      que a função devolve, e o que vale quando \texttt{n} é 0.
    \end{column}
    \begin{column}{0.5\textwidth}
      \small
      \textbf{3. Número mágico}\\
      O \texttt{7.0} é regra da universidade, não da função:
      \texttt{NOTA\_CORTE}, declarada uma vez.

      \medskip
      \textbf{De brinde}\\
      \texttt{const} declara que a função não mexe nas notas --- e
      \texttt{s / n} com \texttt{n} igual a 0 é a pergunta que o comentário
      bom teria levantado.
    \end{column}
  \end{columns}
  \fonte{Kernighan \& Pike (1999), §1.1 e §1.5; McConnell (2004), cap.~11}
\end{frame}

% ==================================================================
\section{O que o compilador cobra}
% ==================================================================

\begin{frame}[fragile]{Número mágico vira constante com nome}
  \begin{columns}[T]
    \begin{column}{0.48\textwidth}
      \textcolor{UFSRed}{\textbf{Código ruim}}
\begin{lstlisting}[style=tinyC]
double R[50];
int Q[50];
...
R[Q[i]] >= 7.0 ? "aprovado"
               : "reprovado"
\end{lstlisting}
      \small O 50 é um limite escondido; o 7.0 é uma regra da universidade
      espalhada pelo código.
    \end{column}
    \begin{column}{0.48\textwidth}
      \textcolor{UFSBlue}{\textbf{No padrão}}
\begin{lstlisting}[style=tinyC]
/* em ranking.h */
#define RANKING_NOTA_CORTE 7.0
\end{lstlisting}
      \small O 50 simplesmente desaparece: o ranking aloca o tamanho da
      turma, como já fazemos desde a aula de alocação dinâmica.
    \end{column}
  \end{columns}
  \fonte{Kernighan \& Pike (1999), §1.5 ``Magic Numbers''}
\end{frame}

\begin{frame}{Variável global: o estado que ninguém vê passar}
  \begin{center}
  \begin{tikzpicture}[
    bx/.style={rectangle, rounded corners, minimum width=2.9cm,
      minimum height=0.8cm, font=\small\ttfamily, align=center}]
    % --- código ruim ---
    \node[bx, draw=UFSRed, fill=UFSRed!8]  (c1) at (0,0)   {calcula()};
    \node[bx, draw=UFSRed, fill=UFSRed!16] (g)  at (4.6,0) {R, Q, N};
    \node[bx, draw=UFSRed, fill=UFSRed!8]  (m1) at (9.2,0) {mostra()};
    \draw[setar] (c1) -- (g)  node[midway, above=1pt, leg] {escreve};
    \draw[setar] (g)  -- (m1) node[midway, above=1pt, leg] {lê};
    \node[leg, anchor=west] at (-1.9,-0.75)
      {a ligação entre as duas não aparece em assinatura nenhuma};

    % --- no padrão ---
    \node[bx, draw=UFSBlue, fill=UFSBlue!8]  (c2) at (0,-2.4)   {ranking\_cria()};
    \node[bx, draw=UFSBlue, fill=UFSBlue!16] (r)  at (4.6,-2.4) {Ranking *r};
    \node[bx, draw=UFSBlue, fill=UFSBlue!8]  (m2) at (9.2,-2.4) {ranking\_imprime()};
    \draw[seta] (c2) -- (r)  node[midway, above=1pt, leg] {devolve};
    \draw[seta] (r)  -- (m2) node[midway, above=1pt, leg] {recebe};
    \node[leg, anchor=west] at (-1.9,-3.15)
      {a ligação está escrita no protótipo --- o compilador cobra};
  \end{tikzpicture}
  \end{center}
  \begin{block}{}
    \small Com globais, chamar \texttt{mostra()} antes de \texttt{calcula()}
    imprime lixo --- e nada no código avisa. Com o parâmetro, não existe
    ranking para imprimir antes de alguém criá-lo.
  \end{block}
  \fonte{Kernighan \& Pike (1999), §4.5 ``Interface Principles''; McConnell (2004), §13.3}
\end{frame}

\begin{frame}[fragile]{\texttt{static}: o que é interno não vaza}
\begin{lstlisting}[style=tinyC]
static void ordena(Ranking *r) { ... }   /* so ranking.c enxerga */
Ranking *ranking_cria(const Turma *t) { ... }  /* interface */
\end{lstlisting}
  \begin{columns}[T]
    \begin{column}{0.5\textwidth}
      \small \textbf{Sem \texttt{static}}\\
      \texttt{troca} e \texttt{maior} entram na tabela de símbolos do
      \texttt{.o}. Outro módulo com uma \texttt{troca} própria derruba a
      ligação com símbolo duplicado.
    \end{column}
    \begin{column}{0.5\textwidth}
      \small \textbf{Com \texttt{static}}\\
      O símbolo não sai do arquivo. A regra fica simples: \emph{está no
      \texttt{.h}? então é público. Não está? então é \texttt{static}.}
    \end{column}
  \end{columns}
  \fonte{ISO/IEC 9899:2018, §6.2.2; Kernighan \& Ritchie (1989), §4.6}
\end{frame}

\begin{frame}[fragile]{\texttt{const}: a promessa de não mexer}
\begin{lstlisting}[style=tinyC]
Ranking *ranking_cria(const Turma *t);   /* le a turma, nao altera */
void ranking_imprime(const Ranking *r);  /* so imprime */
\end{lstlisting}
  \begin{columns}[T]
    \begin{column}{0.5\textwidth}
      \small \textbf{Para quem lê}\\
      O protótipo já responde ``essa função estraga meus dados?'' sem abrir
      o \texttt{.c}.
    \end{column}
    \begin{column}{0.5\textwidth}
      \small \textbf{Para quem escreve}\\
      Se a implementação tentar escrever, o compilador reclama. A promessa
      é verificada, não confiada.
    \end{column}
  \end{columns}
  \begin{block}{}
    \small Já usamos isso sem comentar: \texttt{media(const double *v, int n)}
    e \texttt{turma\_relatorio(const Turma *t)} estão assim desde a aula de
    modularização.
  \end{block}
\end{frame}

\begin{frame}[fragile]{O padrão que a máquina cobra}
\begin{lstlisting}[style=makeStyle]
CFLAGS = -std=c17 -Wall -Wextra -Werror -pedantic -g
\end{lstlisting}
  \begin{columns}[T]
    \begin{column}{0.5\textwidth}
      \small
      \texttt{-Wall -Wextra}: liga os avisos úteis.\\[0.3em]
      \texttt{-Werror}: aviso vira erro --- o projeto não compila com
      pendência.\\[0.3em]
      \texttt{-pedantic -std=c17}: recusa extensões que só o seu compilador
      entende.
    \end{column}
    \begin{column}{0.5\textwidth}
      \small \textbf{Por que já no primeiro dia}\\
      Aviso que ninguém corrige some no meio de outros trinta. Com
      \texttt{-Werror} a dívida nunca acumula: ou o aviso é real e você
      corrige, ou não é e você escreve por que não é.
    \end{column}
  \end{columns}
  \begin{block}{}
    \small E há quem discorde: o GNU Coding Standards diz que
    ``\emph{the compiler should be your servant, not your master}'' e deixa
    o \texttt{-Wall} a critério de cada um. Nenhum dos dois lados está
    errado --- por isso a decisão é da equipe, e a nossa está escrita no
    \texttt{CFLAGS}.
  \end{block}
  \fonte{GCC Manual, ``Options to Request or Suppress Warnings'';
    GNU Coding Standards, §5.3}
\end{frame}

\begin{frame}{Recapitulando: o seu checklist de revisão}
  \begin{columns}[T]
    \begin{column}{0.5\textwidth}
      \small
      $\square$ todo nome público começa com \texttt{ranking\_}?\\[0.35em]
      $\square$ nome de uma letra só dentro de laço curto?\\[0.35em]
      $\square$ nenhuma variável global --- o estado viaja por
        parâmetro?\\[0.35em]
      $\square$ nenhum número solto: o 50 sumiu e o 7.0 tem nome?
    \end{column}
    \begin{column}{0.5\textwidth}
      \small
      $\square$ tudo que não está no \texttt{.h} é \texttt{static}?\\[0.35em]
      $\square$ \texttt{const} em todo parâmetro que a função só lê?\\[0.35em]
      $\square$ cada função do \texttt{.h} diz o que devolve no erro e quem
        libera a memória?\\[0.35em]
      $\square$ o arquivo inteiro tem uma formatação só?
    \end{column}
  \end{columns}
  \begin{block}{}
    \small Cada linha deste checklist é uma decisão que a equipe tomou uma
    vez para não rediscutir a cada arquivo.
  \end{block}
\end{frame}

% ==================================================================
\section{Ferramentas}
% ==================================================================

\begin{frame}[fragile]{Escrever o padrão em um arquivo}
\begin{lstlisting}[style=shellStyle]
# .clang-format --- na raiz do projeto
BasedOnStyle: LLVM
IndentWidth: 4
UseTab: Never
ColumnLimit: 80
BreakBeforeBraces: Linux
AllowShortIfStatementsOnASingleLine: false
\end{lstlisting}
\begin{lstlisting}[style=makeStyle]
formato:
	clang-format -i *.c *.h
\end{lstlisting}
  \begin{block}{}
    \small Com o padrão em arquivo, ninguém precisa lembrar dele --- e a
    revisão de código deixa de gastar tempo com indentação.
  \end{block}
  \fonte{LLVM Project, \emph{ClangFormat Style Options}}
\end{frame}

\begin{frame}{O que a ferramenta faz, e o que ela não faz}
  \begin{columns}[T]
    \begin{column}{0.5\textwidth}
      \textcolor{UFSBlue}{\textbf{Ela resolve}}
      \begin{itemize}
        \item indentação, espaços, chaves;
        \item largura de linha e quebras;
        \item ordem dos \texttt{include}s, se você pedir.
      \end{itemize}
    \end{column}
    \begin{column}{0.5\textwidth}
      \textcolor{UFSRed}{\textbf{Ela não resolve}}
      \begin{itemize}
        \item \texttt{R}, \texttt{Q} e \texttt{N} continuam com esses nomes;
        \item a global continua global;
        \item o módulo continua sem \texttt{.h}.
      \end{itemize}
    \end{column}
  \end{columns}
  \begin{block}{}
    \small \texttt{clang-format} cuida da parte mecânica do padrão. A parte
    que exige entender o programa --- nomes, fronteiras, responsabilidades
    --- continua sendo sua.
  \end{block}
  \fonte{LLVM Project, \emph{ClangFormat}; Kernighan \& Pike (1999), §1.3}
\end{frame}

% ==================================================================
\section{Mãos à obra}
% ==================================================================

\begin{frame}[fragile]{O \texttt{ranking.h} que faltava}
\begin{lstlisting}[style=tinyC]
#ifndef RANKING_H
#define RANKING_H
#include "turma.h"
#define RANKING_NOTA_CORTE 7.0
typedef struct {
    int    *posicao;   /* posicao[0] = indice do aluno de maior media */
    double *media;     /* media[i] = media do aluno de indice i */
    int     quantidade;
} Ranking;

/* Devolve NULL se faltar memoria; quem chama libera. */
Ranking *ranking_cria(const Turma *t);
void     ranking_libera(Ranking *r);
int      ranking_aprovado(const Ranking *r, int colocacao);  /* 0 = 1o */
void     ranking_imprime(const Ranking *r);
#endif /* RANKING_H */
\end{lstlisting}
\end{frame}

\begin{frame}{Altere o programa}
  \small
  \begin{block}{Tarefa 1 \hfill \normalfont\scriptsize agora, em aula}
    Rode \texttt{clang-format -i ranking.c} e leia o \texttt{diff}. Depois
    renomeie \texttt{R}, \texttt{Q}, \texttt{N}, \texttt{calcula} e
    \texttt{mostra}; apague o \texttt{flag}.
  \end{block}
  \begin{block}{Tarefa 2 \hfill \normalfont\scriptsize começa em aula, termina em casa}
    Escreva o \texttt{ranking.h} com o contrato, tire as globais, marque
    \texttt{static} o que não está no header e apague as assinaturas
    copiadas no \texttt{main.c}.
  \end{block}
  \begin{block}{Tarefa 3 \hfill \normalfont\scriptsize em casa}
    Acrescente \texttt{-Werror -pedantic} ao \texttt{CFLAGS} e corrija as
    dependências de \texttt{ranking.o} no \texttt{Makefile}.
  \end{block}
  \begin{center}
    \scriptsize\textbf{Critério de pronto:} compila sem um aviso, \texttt{nm
    ranking.o} só com a interface, saída igual à do começo da aula.
  \end{center}
\end{frame}

\begin{frame}{Desafio}
  \begin{block}{Um critério de ordenação que o cliente escolhe
      \hfill \normalfont\scriptsize em casa}
    Hoje o ranking ordena por média. Faça o critério virar parte da
    interface, sem que \texttt{turma.c} ou \texttt{estat.c} sejam tocados:
    \begin{itemize}
      \item acrescente ao \texttt{ranking.h} um tipo de critério ---
        por média, por maior nota, por menor desvio;
      \item \texttt{ranking\_cria} recebe o critério; a função de comparação
        correspondente fica \texttt{static} dentro de \texttt{ranking.c};
      \item documente no header o que cada critério faz e o que acontece em
        caso de empate;
      \item compile com \texttt{-Werror} e mostre que nenhum símbolo novo
        vazou: confira com \texttt{nm ranking.o}.
    \end{itemize}
  \end{block}
\end{frame}

\begin{frame}{Fechamento}
  \begin{columns}[T]
    \begin{column}{0.5\textwidth}
      \textbf{O que mudou hoje}\\[0.3em]
      O programa imprime exatamente o mesmo que imprimia no começo da aula.
      Mudou quem consegue mexer nele sem medo.
    \end{column}
    \begin{column}{0.5\textwidth}
      \textbf{O fio da aula}\\[0.3em]
      Nome é documentação, formatação é acordo, comentário é o porquê,
      \texttt{static} e \texttt{const} são promessas verificadas e
      \texttt{-Werror} é quem cobra.
    \end{column}
  \end{columns}
  \begin{block}{Para levar}
    \begin{itemize}
      \item padrão de codificação não é gosto: é decisão tomada uma vez para
        não ser rediscutida;
      \item a ferramenta formata; separar responsabilidades continua sendo
        trabalho seu;
      \item código que só você entende é código que só você mantém.
    \end{itemize}
  \end{block}
\end{frame}

\section*{Referências}
\begin{frame}[allowframebreaks]{Referências}
  \nocite{*}
  \bibliographystyle{plain}
  \bibliography{../referencias}
\end{frame}

\end{document}
